%=====================================================================================
\section{Data}\label{sec:data}

\subsection{FOMC statements}

The corpus holds \R{corpus.files} statements released by the FOMC between \R{corpus.first} and
\R{corpus.last}, taken from the press-release pages of the Federal Reserve Board. I classified each by
hand after reading it (Appendix Table~\ref{tab:corpus}): \R{corpus.scheduled} are post-meeting statements
of scheduled meetings, \R{corpus.unscheduledrate} announce unscheduled changes in the federal funds rate,
\R{corpus.unscheduledpolicy} announce unscheduled policy decisions without a rate change (the risk
assessment of 17 August 2007 and the open-ended asset purchases of 23 March 2020), and
\R{corpus.notice} are notices about liquidity facilities, swap lines, technical operations or the
policy framework. The \emph{policy corpus} drops the notices and keeps \R{corpus.policy}
statements. It is used for all reliability statistics and market regressions; the recognition counts in
Section~\ref{sec:recognition} cover every file, notices included. Section~\ref{sec:robust} shows results for a \emph{strict} corpus that
also drops the two unscheduled policy decisions (\R{var.strict.nstmt} statements) and for all
\R{var.all.nstmt} files.

Two corrections were made to the raw corpus. First, the Board's page for the statement of 27--28 June
2007 carries the date 18 June in its address, so the statement was first filed under the wrong day; it
is dated 28 June here, which also gives it its market data. Second, the text of \R{nav.n} statements from
\R{nav.first}--\R{nav.last} originally ended with a line of page navigation that included a year (in
\R{nav.own} cases the statement's own year). That line was removed, but only after the five-prompt scoring
runs had been made, so those runs saw it. The single-prompt re-scores, the GPT-4 run and the dictionary
use the cleaned text (Table~\ref{tab:runs}).

The texts of the statements of 17 June and 29 July 2026 are much shorter than earlier ones and could not
be checked against the Board's website from the computing environment used here. They read as complete
statements in a new, shorter format. They enter the post-2024 checks (Section~\ref{sec:cannot}), the
dating probe on recent statements (Section~\ref{sec:cutoff}) and the comparison of prompt spread on
recent statements (Section~\ref{sec:reliability}); the last is reported with and without them.

\subsection{Market data}

The main outcome is the monetary policy surprise (MPS) of \citet{bauer2023reassessment}. MPS summarises
the change in short-term interest-rate futures in a 30-minute window around each FOMC announcement, in
percentage points; a positive value means markets read the announcement as tighter than expected.
\citet{bauer2023reassessment} show that MPS is partly predictable from data available before the
announcement and also provide MPS\_ORTH, the part of MPS orthogonal to those data. The same dataset
gives the 30-minute changes in the 2- and 10-year Treasury yields ($t_2$, $t_{10}$, in basis points);
their sign and scale match the daily yield changes, with which they correlate at \R{desc.corr.t2.d2}
($t_2$, $d_2$) and \R{desc.corr.t10.d10} ($t_{10}$, $d_{10}$) on statement days. These series end in
December 2023.

Daily yields are the constant-maturity 2- and 10-year Treasury yields from the Federal Reserve's H.15
release (FRED series DGS2 and DGS10). The daily outcomes $d_2$ and $d_{10}$ are close-to-close changes in
basis points on the statement day, from the previous trading day. The yield file matches the FRED series
to the hundredth on all \R{fred.days} days on which they overlap (\R{fred.first} to \R{fred.last}). A
statement released on a non-trading day (15 March 2020, a Sunday) takes the next trading day's change.
The daily series used here starts in 2003, so the daily regressions have fewer meetings than those using
MPS.

A \emph{hold meeting} is a statement at which the federal funds rate target did not change. Hold meetings
matter for this paper because at them the stance score cannot simply echo a rate move. They are not free
of policy news: a hold can itself surprise markets, and the hold sample includes announcements of asset
purchases (for example 18 March 2009) and the unscheduled statements of 17 August 2007 and 23 March
2020. Before 2024 the policy
corpus contains \R{desc.pre} statements, which give \R{desc.prepairs} meeting-to-meeting changes;
\R{desc.hold} of these are hold meetings with an MPS. Table~\ref{tab:desc} describes the outcomes.

\begin{table}[t]
\centering
\caption{Market outcomes, policy corpus, statements to December 2023}\label{tab:desc}
\small
% generated by paper/make_tables.py - do not edit
\begin{tabular}{@{}llrrrrrrl@{}}
\toprule
Outcome & Unit & Sample & $n$ & Mean & SD & Min & Max & Years \\
\midrule
MPS & pp & all & 198 & $-$0.002 & 0.059 & $-$0.283 & 0.152 & 2000--2023 \\
 &  & hold & 132 & 0.002 & 0.039 & $-$0.111 & 0.152 & 2000--2023 \\
MPS\_ORTH & pp & all & 189 & 0.003 & 0.052 & $-$0.232 & 0.161 & 2000--2023 \\
 &  & hold & 125 & 0.003 & 0.041 & $-$0.107 & 0.161 & 2000--2023 \\
2-year yield, 30-minute ($t_2$) & bp & all & 198 & $-$0.8 & 5.0 & $-$16.9 & 20.5 & 2000--2023 \\
 &  & hold & 132 & $-$0.7 & 4.4 & $-$13.7 & 20.5 & 2000--2023 \\
10-year yield, 30-minute ($t_{10}$) & bp & all & 198 & $-$0.3 & 4.2 & $-$29.6 & 16.4 & 2000--2023 \\
 &  & hold & 132 & $-$0.4 & 4.4 & $-$29.6 & 16.4 & 2000--2023 \\
2-year yield, daily ($d_2$) & bp & all & 173 & $-$0.9 & 7.7 & $-$28.0 & 23.0 & 2003--2023 \\
 &  & hold & 120 & $-$0.8 & 6.1 & $-$27.0 & 17.0 & 2003--2023 \\
10-year yield, daily ($d_{10}$) & bp & all & 173 & $-$0.9 & 8.0 & $-$51.0 & 22.0 & 2003--2023 \\
 &  & hold & 120 & $-$0.5 & 8.0 & $-$51.0 & 20.0 & 2003--2023 \\
\bottomrule
\end{tabular}

\tabnote{\emph{Notes:} ``all'' is every statement with a predecessor in the corpus; ``hold'' is the
subset at which the federal funds rate target did not change. MPS and MPS\_ORTH are from
\citet{bauer2023reassessment}, in percentage points. $t_2$ and $t_{10}$ are 30-minute changes in the 2-
and 10-year Treasury yields around the announcement, from the same dataset; $d_2$ and $d_{10}$ are daily
close-to-close changes of the H.15 constant-maturity yields (FRED DGS2, DGS10).}
\end{table}

%=====================================================================================
\section{Scoring design}\label{sec:design}

\subsection{Models and prompts}

Every statement is scored by asking a model for a single number between $-1$ (dovish) and $+1$
(hawkish). I use five prompts that ask the same question in different words (Appendix~\ref{app:prompts}
gives them verbatim). P1 asks for the score with no further instruction. P2 first defines hawkish and
dovish. P3 asks the model to answer as a fixed-income analyst. P4 reverses the scale, so that $+1$ means
dovish; its answer is flipped back in code. P5 asks the model to list the hawkish and dovish elements
before giving the score. All calls use temperature 0 through the OpenRouter interface.

Table~\ref{tab:runs} lists the scoring runs. The reliability block crosses \R{icc.fable.n} policy
statements from 2015--2019 with all five prompts and four models: Gemini 2.5 Flash, Claude 3 Haiku,
GPT-4o mini and Claude Fable 5.1. Gemini and Fable also score the full corpus. Gemini has all five
prompts for every policy statement. Fable has all five for \R{cover.fable.five} statements and between
one and four for the other \R{cover.fable.fewer}, all of them from \R{cover.fable.fewfirst} onwards. GPT-4, listed by OpenRouter
with a training cutoff of 30 September 2021, scores every statement with P1. Fable then re-scores the
pre-2024 statements with P1 three times: on the cleaned original text, on a blinded version and on a
paraphrased version (Section~\ref{sec:versions}).

Replies are parsed strictly: the reply must end in a number in $[-1,1]$, or contain a line
``SCORE:~$x$''. A reply that is cut off, ends in a bare ``0.'' or gives a value outside the range is
treated as missing. Until 3 October 2026 the scoring scripts stored only the first 300 characters of each
reply, so longer replies could not be re-parsed afterwards. Re-parsing every shorter stored reply and
reading one longer reply by hand removed \R{parse.dropped} cells (Appendix Table~\ref{tab:parse}).

\begin{table}[t]
\centering
\caption{Scoring runs}\label{tab:runs}
\footnotesize
\setlength{\tabcolsep}{4pt}
% generated by paper/make_tables.py - do not edit
\begin{tabular}{@{}llrrl@{}}
\toprule
Model (OpenRouter ID) & Use & Prompts & Statements & Text \\
\midrule
\texttt{google/gemini-2.5-flash} & Reliability block & 5 & 41 & original \\
\texttt{anthropic/claude-3-haiku} & Reliability block & 5 & 41 & original \\
\texttt{openai/gpt-4o-mini} & Reliability block & 5 & 41 & original \\
\texttt{anthropic/claude-fable-5.1} & Reliability block & 5 & 41 & original \\
\texttt{google/gemini-2.5-flash} & Reliability, 2004--07 & 5 & 34 & original \\
\texttt{google/gemini-2.5-flash} & Full corpus & 5 & 228 & original \\
\texttt{anthropic/claude-fable-5.1} & Full corpus & 5 & 228 & original \\
\texttt{openai/gpt-4} & Full corpus (cutoff model) & 1 & 228 & original, cleaned \\
\texttt{anthropic/claude-fable-5.1} & Re-score, pre-2024 & 1 & 206 & original, cleaned \\
\texttt{anthropic/claude-fable-5.1} & Re-score, pre-2024 & 1 & 206 & blinded \\
\texttt{anthropic/claude-fable-5.1} & Re-score, pre-2024 & 1 & 206 & paraphrased \\
word count (no model) & Full corpus (no LLM) & 1 & 228 & original, cleaned \\
\bottomrule
\end{tabular}

\tabnote{\emph{Notes:} Model IDs as listed by OpenRouter. Statements = distinct statement files with a score (all 228 files, including the
notices that the policy corpus drops). The full-corpus Fable file has all five prompts for
\R{cover.fable.five} policy statements and one to four for \R{cover.fable.fewer}. ``Original'' text is the text as first downloaded; for \R{nav.n}
statements from \R{nav.first}--\R{nav.last} it ended with a navigation line that included a year.
``Cleaned'' text has that line removed. Blinded and paraphrased versions are described in
Section~\ref{sec:versions}. The dictionary is a word-count score (Section~\ref{sec:baselines}).}
\end{table}

\subsection{From scores to dScore}

For each statement and prompt I average repeated runs of the same cell; Fable's test--retest agreement on
the \R{retest.n} cells scored twice is high (correlation \R{retest.r}, mean absolute difference
\R{retest.mad}). The statement's score is the mean over its prompts after removing each prompt's average
offset, estimated on the statements that have all five prompts (prompt fixed effects). Without this
step, a statement scored with a different set of prompts than its predecessor would show a spurious
change. The score change, \emph{dScore}, is the score of a statement minus the score of the previous
statement in the policy corpus. Notices are removed before the differences are taken, so each statement
is compared with the previous policy statement. A positive dScore means the statement reads as more
hawkish than its predecessor. Section~\ref{sec:robust} shows that the main results do not depend on
the prompt fixed effects.

\subsection{Text versions}\label{sec:versions}

The \emph{blinded} version of each statement removes the cues that most directly date it. Month names
become ``[month]''; years become offsets from the statement's own year (``[year+1]''), so that date-based
guidance keeps its meaning; the names of all FOMC members become ``[member]''; the roster of members
voting in favour is dropped, while dissents are kept without names; and named one-off events become
``[event]''. Rate levels, purchase amounts and the phrasing are left as they are.

The \emph{paraphrased} version rewrites each blinded statement with Gemini 2.5 Flash. The instructions
(Appendix~\ref{app:prompts}) ask for the same economic content in a new order and new words, with no
reused phrase of four or more words, and with rate levels and asset holdings stated only as changes from
the previous meeting. To keep the decision right, the prompt states the true change in the federal funds
target. Automatic checks flag rewrites that are too long, add dissents or a 2 percent objective that the
original lacks, or keep rate levels; a flagged rewrite is tried once more. Of the \R{para.n}
paraphrases, \R{para.passed} pass the checks and \R{para.flagged} (\R{para.flagged.dates}) still keep a
number after the retry: a guidance threshold in the first case, an operational rate in the second, which
is a notice outside the policy corpus. A retry was needed for \R{para.retries} rewrites. Their four-word overlap with the blinded text is
\R{para.overlap.mean} on average (maximum \R{para.overlap.max}), and they are \R{para.length} times as
long. In all, \R{para.sixhalf} paraphrases keep the 6.5 percent unemployment-rate threshold from the
forward guidance of 2012--2014; the checks did not treat a threshold as a level. I read a sample of
paraphrases against their originals, including those on which an earlier attempt had gone wrong, and
found the content faithful.

\subsection{Baselines}\label{sec:baselines}

Two baselines involve no LLM. The \emph{dictionary} score counts hawkish and dovish word combinations
(for example, ``inflation'' near ``rise'' counts as hawkish) in the spirit of
\citet{apel2012}, using a short word list written for this project, and scales the net count to
$[-1,1]$. It cannot have memorised anything. The \emph{text-similarity} dating baseline, used in
Section~\ref{sec:recognition}, gives each blinded statement the date of the most similar other blinded
statement. Finally, \citet{shah2023} hand-label sentences from FOMC minutes, speeches and press
conferences as hawkish, dovish or neutral. Because the minutes repeat the statement, \R{tdw.matches} of
their labelled sentences appear word for word in \R{tdw.stmts} of the statements here. This gives a human
benchmark for score levels that neither markets nor model memory can affect.
